\cotefile{78}
\scanpage{61}

Pour prouver \unsure{$\Leftarrow$)|$\Leftarrow$|a)}, notons que $\sigma_{i}R^+=R^-$ $\Longleftrightarrow \sigma_i R^-=R^+\unsure{,|{}}$ $(i\in\{0,\unsure{2|1}\})$, donc si $R^+$ est une orientation, on aura $uR^+=R^+$ donc $G_1^+(R^+)\subset R^+$, et $gR^+=R^-$ si $g\in G_1\setminus G_1^0=G_1^-$, \deleted{Il suffit de voir que} et inversement. Or dans le cas d'un polygone, \deleted{\unsure{$\simeq G$}} \unsure{donc} $\simeq\Pi_{n}$, on voit immédiatement que $G_1^+$ opérant sur $R\simeq\mathbb D_n$ a \inserted{\unsure{exactemt|exactement}} deux orbites, \deleted{qui} \inserted{lesquelles sont} échangées par $\sigma$, savoir $D_n^+$ et $D_n^-$, -- ce sont donc \unsure{celà|là|ici} les orientations.

\deleted{\illegible{1}} \textcircled{7}. \uline{Prop} Soit $\Pi$ \inserted{$=(S,A,R)$} polygone combinatoire, \deleted{\unsure{$G$}} \inserted{désignons} $G=\operatorname{Aut}\Pi$, \deleted{\unsure{On|donc}} $\operatorname{Or}(\Pi)$ a deux éléments et \deleted{\unsure{pour|par|$G$}} transport de structure \dots $G\longrightarrow\unsure{\mathfrak S|\mathscr S}_{\operatorname{Or}(\Pi)}\simeq\pm1$, \deleted{\unsure{Rappelons|que l'on}} Soit $G^+$ le noyau, a) On \deleted{\unsure{obt}} \inserted{obtient ainsi} suite exacte

\[
1\longrightarrow G^+\longrightarrow G\mathbin{-}\{\pm1\}\mathbin{-}1
\]

et tt élément de $G^-=G\setminus G^+$ est d'ordre $2$, donc splitte cette suite exacte. \deleted{b) $R$ est torseur} \deleted{\unsure{a)}} $G^+$ est cyclique d'ordre $n$, \deleted{si $R^+\in\operatorname{Or}(\Pi)$, et $S,A$ \unsure{sont des}} \inserted{et l'él. $-1$ de $\pm1$} y opère par symétrie.

b) $R$ est un torseur sous $G$, $S,A$ sont torseurs sous $G^+$.

% archive sheet number: 60
\scanpage{62}

c) Pour tt $\omega\in\operatorname{Or}(\Pi)$, $\exists\unsure{1|!}u_\omega\in G^+$ tel que $u_\omega\mid R_\omega^+=u\mid R_\omega^+$, et on a \deleted{\unsure{$u_\omega$}}

\[u_{-\omega}=u_\omega^{-1}\]

\deleted{Donc} De plus, $u_\omega$ est un générateur de $G^+$, et si $n\geq 3$, $\omega\mapsto u_\omega$ est une injection.

\[\operatorname{Or}(\Pi)\unsure{\leftrightarrow|\hookrightarrow}\operatorname{Génér.}(G^+)\]

\margin{\deleted{\unsure{$\nearrow$}}\ $\nearrow$\ $\nearrow$}
% slanted

d) Le foncteur

\[\Pi\longmapsto (Z=G^+,\inserted{\Omega=}\{u_\omega,u_{-\omega}\}\in G^+,S)\]

des $n$-polygones \inserted{$(\unsure{3|{-3}}\leq{n}\leq+\infty)$} vers les triples d'un groupe cyclique d'ordre $n$ $G^+$, d'un ens de deux gén. de $G^+$ inv. l'un de l'autre, et d'un $G^+$ torseur $S$, est une équivalence de \unsure{ctégories|catégories}.

\uline{\unsure{Exer|Exo|Exer.}} Exprimer $A$ \inserted{(en \unsure{tt|tant} que $Z$-torseur)} en termes des données précédentes, ainsi que $R$ en tt que partie de $S\times R$ \dots

% archive sheet number: 61
\scanpage{63}

\begin{tabular}{|l}Cas des $\unsure{n|m}$-polygônes réguliers, \deleted{\unsure{de}} $n$ premier à la\\ caractéristique ($k$ alg clos)\end{tabular}

$S$ ensemble des sommets, $A$ ens des arêtes, $E$ espace affine ambiant.

Le barycentre $o$ de $S$ est fixé par le groupe $G$ des automorphismes de $\Pi$, de sorte que $E$ devient un vectoriel. Soit $G^+$ le sous-groupe des ``rotations'' de $\Pi$, $u$ un des $2$ générateurs privilégiés, correspondant à une orientation \inserted{combinatoire} de $\Pi$, \deleted{Considérons} \inserted{soit $s_0\in S$} \deleted{les deux valeurs propres de $u$, \unsure{puis}} [la donnée de $(s_0,u)$ équivaut à celle d'un repère combinatoire, i.e. d'un épinglage combinatoire] $\sigma$ l'autom $\ne\mathrm{id}$ tel que $\sigma s_0=s_0$, de sorte que

\[\left\{\begin{aligned}
&\sigma^2=\mathrm{id},\quad \sigma u\sigma^{\unsure{-1|\prime\prime}}=u^{-1}\\
&u^n=1\qquad\text{(pour mémoire)}
\end{aligned}\right.\]

Soient $\zeta,\zeta'$ les valeurs propres de $u$. Si on avait $\zeta=\zeta'$, on aurait $u=\zeta(\mathrm{id}_E+v)$, $v$ nilpotent, \inserted{i.e. $v^2=0$} et $u^n=\mathrm{id}\Longrightarrow\zeta^n=1$, $1+nv=0$ i.e. $nv=0$, d'où (comme $n1_{\unsure{k|K}}\ne0$ par hyp) $v=0$. donc $u=\zeta\,\mathrm{id}_E$, donc les $s_i=u^is_0$ seraient sur une $\overline{\mathrm m}$ droite, \uline{absurde}. Donc $\zeta\ne\zeta'$, soient $D_\zeta,D_{\zeta'}$ les droites propres. \deleted{\unsure{Mais}} Comme $u\sim u^{-1}$ l'un des

% archive sheet number: 62
\scanpage{64}

valeurs propres est stable par $z\mapsto z^{-1}$, i.e \deleted{\unsure{$\zeta,\zeta'$|$\zeta\zeta'$}}

\[\zeta\zeta'=1\qquad\text{ou}\qquad\zeta'=\zeta^{-1}\]

plus précisément, comme $\sigma u\sigma^{\unsure{-1|\prime\prime}}=u^{-1}$ $\sigma$ échange $D_\zeta$ et $D_{\zeta'}$. \deleted{Choisissons $e_1\in D_\zeta$ base, d'où} \deleted{$e_2=\sigma e_1$ base de \unsure{$D_{\zeta'}$}, et} \inserted{Si $s_0=e_1+e_2$ ($e_1\in D_\zeta$, $e_2\in D_{\zeta'}$) \deleted{\unsure{de $D_{\zeta'}$}} alors $s_0=\sigma s_0=\sigma e_2+\sigma e_1$ donc $\sigma e_1=e_2$, $\sigma e_2=e_1$ -- on a aussi $e_1,e_2\ne0$ (sinon on aurait $s_0=0$, donc $s_i=s_0$ tous, absurde) $e_1,e_2$ base de} $E$ (définie \deleted{\unsure{modulo multiplication par scalaire}} \inserted{\unsure{canoniquement|canon\textsuperscript{t}|de façon unique}, \inserted{modulo permutation $e_1\leftrightarrow e_2$,} \unsure{ne dépend que de|en termes de} l'épinglage}).

Dans cette base, on a

\[\begin{aligned}
u&=\begin{pmatrix}\zeta&0\\0&\zeta^{-1}\end{pmatrix}&\sigma&=\begin{pmatrix}0&1\\1&0\end{pmatrix}\\
s_0&=\begin{pmatrix}1\\1\end{pmatrix}&s_i&=\begin{pmatrix}\zeta^i\\\zeta^{-i}\end{pmatrix}
\end{aligned}\]

On est dans le cas-typique déjà étudié (sans avoir eu à supposer précédemment $k$ alg. clos, ni $\zeta$ d'ordre fini \dots). De plus, on voit que

\deleted{\uline{Fonction quadratique invariante} (On avait déjà remarqué que (p.r. à la base envisagée)}

\[\alpha=\zeta+\zeta'=\zeta+\zeta^{-1}=\operatorname{Tr}u=\operatorname{Tr}u^{-1}\]

est défini de façon intrinsèque en termes de $\Pi$, ou ce qui revient au $=$, la \deleted{\unsure{couple|couplée} \unsure{$(\zeta,$}} \inserted{\uline{paire} (non ordonnée)} $\{\zeta,\zeta'\}=\{\zeta,\zeta^{-1}\}=\{\zeta',{\zeta'}^{-1}\}$ d'éléments de $k^*$ inverses l'un de l'autre est défini intrinsèquement en termes de $\Pi$

% archive sheet number: 63
\scanpage{65}

On peut prendre pour $\zeta$ n'importe quelle racine de $1$, à l'exclusion de $\pm1$ (racines d'ordre $1$ et $2$ \dots).

Ainsi on a trouvé

\uline{Proposition} a) Tout \deleted{\unsure{épinglé|choix}} polygône régulier épinglé à $n$ cotés, $3\leq n<+\infty$ (donc $n$ fini), sur un corps alg. clos $k$ de car. $p$ premier à $n$, est isomorphe à un \inserted{des} polygônes épinglés $(\Pi_{\zeta,k},\deleted{\unsure{s}}\,\unsure{r_{n,k}|r_{\zeta,k}})$ \inserted{dans $k^2$} construits précédemment, \inserted{où $\zeta$ est racine primitive $n$ ième de $1$.} Pour \unsure{$\Pi,r$|$(\Pi,r)$|$\Pi,n$|$\Pi$} fixé, $\zeta$ est défini modulo passage à $\zeta^{-1}$, \deleted{\unsure{on}}

b) On trouve ainsi une correspondance $1$-$1$ entre classes d'isom. de $n$-polygônes réguliers sur $k$ (ou classes d'isom de $n$-polygônes réguliers épinglés), entre \deleted{couples} \inserted{paires} $\{\zeta,\zeta^{-1}\}$ de paires de racines \unsure{primes|primièmes|primitives} $n$\textsuperscript{ièmes} de $1$ inverses l'une de l'autre, et entre éléments $\alpha\in k$ tels que les \inserted{$2$} solutions de l'équation

\margin{terminologie: \uline{invariant} \uline{numérique} des polygônes réguliers \dots}
% slanted

\[\zeta^2\deleted{\unsure{+}}-\alpha\zeta+1=0\]

soient des racines primitives $n$\textsuperscript{ièmes} de $1$.

c) Dans $\Pi_{\zeta,k}$, l'origine de $k^2$ est \unsure{car.|aussi} comme le barycentre de $S$, ou encore l'unique pt fixe sous $G$\inserted{$=\operatorname{Aut}(\Pi_{\unsure{n|m|\zeta},k})$}, ou l'unique pt fixe sous $G^+$. $\{\zeta,\zeta^{-1}\}$ est l'ens des valeurs

% archive sheet number: 64
\scanpage{66}

propres de la partie homogène de \deleted{\unsure{$u$ ou $u^{-1}$}} \inserted{l'un} quelconque ($u$ ou $u^{-1}$) des deux gén privilégiés de $G^+$, les axes de coordonnées sont les droites propres \deleted{pour} \inserted{de} $u$ (associé à l'épinglage) pour les val. propres $\zeta$ et $\zeta^{-1}$ resp\textsuperscript{t}.

\begin{center}\rule{1.5cm}{0.4pt}\end{center}

\uline{Formes} \deleted{\uline{quadratiques}} \uline{invariantes} (sous $\operatorname{Aut}(\Pi_{\zeta,k})$)

Si on regarde les fonctions \unsure{polygôniales|polygônales|polynômiales} \deleted{\illegible{1}} \inserted{sur} $E$ invariantes par $G$, comme l'origine est fixe, elles se décomposent en fonctions polygôniales homogènes, \deleted{En degré $m$} ou \uline{formes}

L'invariance par $\sigma$ signifie la symétrie p.r. aux \deleted{$x,y$} coordonnées $x,y$, donc

\[\left\{\begin{aligned}
&F(x,y)=G(\sigma_1(x,y),\sigma_2(x,y))\\
&\sigma_1(x,y)=x+y,\quad \sigma_2(x\unsure{.|,}y)=xy
\end{aligned}\right.\]

D'ailleurs, si \unsure{;|,|{}}

\[\deleted{\unsure{\operatorname{ord}}}F(x,y)=\sum_{\alpha+\beta=m}c_{\alpha\beta}\,x^\alpha y^\beta\]

on a

\[F(u(x,y))=F(\zeta x,\zeta^{-1}y)=\sum c_{\alpha\beta}\,\zeta^{\deleted{\unsure{\beta}}\alpha-\beta}x^\alpha y^{\unsure{\beta|\beta\zeta}}\]

donc on a

\[F\equiv F\circ u\Longleftrightarrow c_{\alpha\beta}=0\quad\text{si}\quad\beta-\alpha\not\equiv0\ (n)\]

% archive sheet number: 65
\scanpage{67}

Donc les \deleted{\unsure{poly}} fonctions \unsure{polygôniales|polynômiales} invariantes ont une base sur $k$ formée de $1$ et des

\[\begin{aligned}
&\deleted{\varphi_\alpha(x,y)}\\
\varphi_{\alpha,i}(x,y)&=x^\alpha y^{\alpha+in}+y^\alpha x^{\alpha+in}\\
&=x^\alpha y^\alpha(x^{in}+y^{in})\\
&=(xy)^\alpha\big(\underbrace{(x^n)^i+(y^n)^i}\big)
\end{aligned}\]

avec $\alpha\in\mathbb N$, $i\in\mathbb N^*$

\deleted{\unsure{Soit}} Or si $X,Y$ sont des indéterminées, le Th. des fonctions sym. nous dit que

\[\begin{aligned}
X^i+Y^i&=\underset{\inserted{\substack{\searrow\ \text{\unsure{polygône|polynôme} universel à coeff}\ \in\mathbb Z\\
\text{en }2\text{ variables }\sigma_1,\sigma_2}}}{S_i}(\sigma_1(X,Y),\sigma_2(X,Y))\\
&=S_i(X+Y,XY)
\end{aligned}\]

donc

\[(x^n)^i+(y^n)^i=S_i(x^n+y^n,(xy)^n)\]

et par suite

\[k[x,y]^{\overset{\inserted{\nearrow\,=\operatorname{Aut}(\Pi_{\zeta,k})}}{D_{\unsure{\zeta|G|{}}\deleted{\unsure{,k|\zeta}}}}}
=k[x,y]^{\overset{\inserted{\substack{\nearrow\,\text{groupe }\mu_n\mathbin{.}\mathbb Z/2\mathbb Z\ (\unsure{\tfrac12\,\text{direct}|\tfrac12\,\text{dirct}})\\
\text{opérant sur }k^2\text{ de façon évidente}}}}{\Delta_n}}=k[\sigma_2,S_n]\]

où

\[\left\{\begin{aligned}
&\sigma_2(x,y)=xy\\
&S_n(x,y)=x^n+y^n
\end{aligned}\right.\]

\uline{NB} Les calculs marchent pour tt $n\in\mathbb N^*$ (pas besoin que $n\geq3$) et sur tt anneau de base -- le meilleur de tous étant $\mathbb Z$ ! Si $n$ est infini, il faut revenir au cas d'un $\zeta$

% archive sheet number: 66
\scanpage{68}

donné, section d'un faisceau structural, d'ordre infini sur chaque fibre, on trouve comme anneau d'invariants $k[\sigma_2]$ \dots.

\uline{Corollaire} Les fonctions polynômiales de degré $\leq2$ invariantes par $D_\zeta=\operatorname{Aut}(\Pi_{\zeta,k})$ sont les fonctions

\[axy+b\qquad (a,b\in k)\]

Donc à scalaire près une unique forme quadratique \unsure{invar.|invariante}, \inserted{$q(x,y)=\deleted{\unsure{q}}\sigma_2(x,y)=xy$} qu'on peut normaliser par $\sigma_2(\unsure{S|s})=1$ si $s\in S$. \uline{Une unique} \uline{cônique affine} ``circonscrite à $\Pi$'' i.e. contenant les sommets, savoir celle d'équation

\margin{\unsure{vif !|vérif !|réf !|Kif !} Ne marche que si $n\geq5$, avec démonstration à donner. Ici on a néanmoins une cônique circonscrite canonique -- invariante par $G$}
% slanted

\[xy-1=0\qquad\deleted{\unsure{xy-1=0|xy-1=2}}\]

\deleted{La forme quadratique $\sigma_2$}

Pour un $n$-polygône régulier quelconque (sur $k$ alg. clos de car première à $n$) -- pas épinglé -- il y a donc \inserted{à scalaire près} une unique forme quadratique invariante $\unsure{\zeta_0|q_0}$, qu'on normalise par

\[q_{0}(s)=1\qquad(s\in S)\]

% archive sheet number: 67
\scanpage{69}

et une unique \deleted{\unsure{quadrique|quadratique}} \inserted{conique} circonscrite à $\Pi$ i.e. à $S$, savoir la cônique

\[C\ \text{d'équation}\qquad q_{0}(x)-1=0\]

Notons que $q_0$, considéré comme forme quadratique sur l'espace des translations de $E$, fait de $E$ un ``plan \uline{euclidien}'' -- la forme est en effet non dégénérée. Le carré de la longueur des cotés du polygône $\Pi$ est donné par

\[\begin{aligned}
q_0(s_1-s_0)&=q_0(\zeta-1,\zeta^{-1}-1)=(\zeta-1)(\zeta^{-1}-1)\\
&=2-\alpha
\end{aligned}\]

i.e.

\[q_0(a)=2-\alpha\]

On notera que c'est tjrs $\ne0$ i.e. \deleted{\unsure{$2-\alpha$}} $\alpha\ne2$, puisque $\zeta\ne1$. [On pourrait donc normaliser \deleted{\unsure{$q_0$|$q$}} $q$ \unsure{différemt|différemment}, en normalisant par $q_0(a)=1$. Dans le cas général, c'est cette normalisation-là qui sera la meilleure, comme nous verrons \dots.]

\deleted{\uline{Commentaires ``philosophiques''} On a trouvé, pour le gpe $D_{n}$ fixé, et des corps}

% archive sheet number: 68
\scanpage{70}

\deleted{\deleted{(alg clos} variables $k$ \inserted{alg clos (pour simplifier)} de car. première à $n$, \deleted{des} \uline{munis} d'un $\alpha\in k$ tel que les solutions de $\zeta^2-\alpha\zeta+1=0$ soient des racines primitives $n$\textsuperscript{ièmes} de $1$ (ce qu'on peut exprimer par la condition $F_n(\zeta)=0$ \dots) une \deleted{\unsure{rep}} famille de $k$-représentations affines de dim $2$ \inserted{de $D_n$} \unsure{ttes|tous|des} isomorphes -- savoir celles associées aux \deleted{\inserted{représentations} polygônes réguliers} couples formés d'un polygône régulier d'invariant $\alpha$, et d'un épinglage de celui-ci -- ou encore des hom $D_n\longrightarrow\operatorname{Aff}\unsure{(2k)|(2,k)|(3k)}$ définis à autom intérieur près.}

\uline{\unsure{Remarques|Remarque}}

Si on s'était borné d'exiger, pour un polyèdre \inserted{affine} \uline{régulier}\deleted{affine}, que le gpe de ses autom \uline{affines} soit transitif sur l'ens des sommets, \inserted{(sans \unsure{épuiser|préciser|exiger} \unsure{néc|néc\textsuperscript{t}|rien} \unsure{l'ens|leurs} \deleted{\unsure{groupe}} des autom \unsure{comb|combs|ensembles})} on aurait trouvé beaucoup plus de \unsure{classes|choses}. Par exemple, prenant le cas particulier où \inserted{le gpe $G^+$} des rotations combinatoires provient d'un gpe de transformations affines, on aurait eu à classer des $u$ tels que $u^n=1$ (sans disposer de $\sigma$ tel que $\sigma u\sigma^{-1}=u^{-1}$) ce qui nous aurait donné deux val. propres $\zeta,\zeta'$ racines \deleted{\unsure{pri|pr}} $n$\textsuperscript{ièmes} de $1$, d'ordre $d,d'\mid n$, avec

% archive sheet number: 69
\scanpage{71}

$n=\operatorname{ppcm}(d,d')$ (pour exprimer que $u$ est d'ordre $n$ exactement) -- $S$ se définissant en prenant l'orbite d'un $s_0\in E\setminus(D_\zeta\cup D_{\zeta'})$, \deleted{et $A$} i.e. l'ens des $s_i=u^i s_0$, \deleted{et} $(i\in\mathbb Z/n\mathbb Z)$, et $A$ est formé des droites $a_i=\operatorname{dr}\{s_i,s_{i+1}\}$ (il faut encore vérifier, ou exiger, que $a_i$ n'est incident qu'à $s_0$ et $s_1$, cette condition est laissée en exercice...

\rule{1.7cm}{.35pt}

Par contre, si on exige \inserted{de plus} qu'il existe une \deleted{\unsure{s|$\delta$}}métrique non dégénérée invariante par les $\unsure{s|u|\delta}\in G^+$, et en insistant pour que les automorphismes soient des \uline{rotations} (comme l'avait suggéré \unsure{Olivier Laborde|Olivier Leborgne}), on sauve la mise. En effet, les rotations sont unimodulaires, donc on trouve que les valeurs propres de $u$ sont $\{\zeta,\zeta^{-1}\}$, et on retrouve les cas précédents.... Je n'ai pas vérifié si ça marche encore en exigeant simplement que les $\in G^+$ respectent la métrique, sans plus.....

\rule{2.2cm}{.35pt}

% archive sheet number: 70
\scanpage{72}

Fonctions polynômiales rel\textsuperscript{t} invariantes sous $G^+$, sous $G$.

\deleted{\unsure{1}}. Les $x^{i}y^{\unsure{mj|j}}$ forment une base de $\operatorname{Sym}^{\unsure{*|\bullet}}(\check E)$, formée de vecteurs propres sous $u$, de valeur propre $\zeta^{\unsure{i-mj|i-j}}$. \deleted{\illegible{1}} Les seules valeurs propres de $G^+\simeq\mathbb Z/n\mathbb Z$ opérant sur $\operatorname{Sym}^*(\check E)\simeq\text{Fonctions polyn. }(E)$\deleted{\unsure{$\ne$|$k$}} sont donc les $\zeta^d$ \inserted{$(d\in\mathbb Z/n\mathbb Z)$}, et les fonctions correspondant à cette valeur propre sont les combinaisons linéaires des \deleted{$x^{j+d}y^j=(xy)^j\unsure{x|x^d}$} i.e.

\[x^iy^j\qquad i-j\equiv d\quad(n)\]

i.e. les

\[
x^{j+d\deleted{\unsure{k}}+nk}y^j=
\left\{\begin{aligned}
&(xy)^j x^{d+nk}&&\text{si }d+nk\geq0\\
&(xy)^{\unsure{i|j|j'}}y^{-d+n\deleted{\unsure{+}}k'}&&\text{si }-d+nk'\geq0
\end{aligned}\right.
\]

\deleted{\unsure{En tant d|Cas $d=0$}} En tant que module sur le sous anneau

\[\operatorname{Sym}^*(\check E)^{G^+}\simeq k[
\overset{\inserted{\substack{W=W_n\\ \Vert}}}{xy},
\overset{\inserted{\substack{U=U_n\\ \Vert}}}{x^n},
\overset{\inserted{\substack{V=V_n\\ \Vert}}}{y^n}]
\quad(U_{n}V_n-W_n^n=0)\]

c'est engendré \deleted{\unsure{si d}} \deleted{par $x^d$ et $y^{-d+\unsure{n|2}k}$}, ou si $0<d<n$, par $\underset{\inserted{\substack{\Vert\\X_d}}}{x^d}$ et $\underset{\inserted{\substack{\Vert\\Y_d}}}{y^{n-d}}$, avec les conditions $VX_d-(W)^dY_d=0$ ....

Pour les fonctions polynômiales de degré $\leq2$,

% archive sheet number: 71
\scanpage{73}

ayant donc la base

\[1,x,y,x^2,xy,y^2\]

les valeurs propres de $u$ correspondantes sont resp\textsuperscript{t}.

\[1,\zeta,\zeta^{-1},\unsure{\zeta|\zeta^2},1,\zeta^{-2}\]

donc les valeurs propres qui y figurent sont \deleted{\unsure{respectivement}} parmi les \inserted{$5$} valeurs

\[\zeta^{-2},\zeta^{-1},\zeta^0=1,\zeta^1=\zeta,\zeta^2.\]

Celles-ci sont \deleted{\unsure{ttes|toutes}} distinctes \unsure{ssi|sss|si} $n\geq5$, dans ce cas les \deleted{fonctions}\inserted{poly} quasi-invariantes de degré $\leq2$ \inserted{sous $G^+$} sont exact\textsuperscript{t} les fonctions

\[\begin{aligned}
a+bxy&&\text{valeur propre }&\zeta^0=1\\
ax\quad(\text{resp. }ay)&&\text{valeur propre }&\zeta\quad(\text{resp }\zeta^{-1})\\
ax^2\quad(\text{resp }ay^2)&&\text{valeur propre }&\zeta^2\quad(\text{resp }\zeta^{\unsure{-1|-2}})
\end{aligned}\]

Si on veut de plus la quasi-invariance sous $G\simeq D_n$ tout entier, \deleted{sachant que $\sigma$}\inserted{i.e. sous $\sigma$} notant que $\sigma(ax)=ay$, $\sigma(ay)=ax$, $\sigma(ax^2)=ay^{2}$, $\sigma(ay^2)=ax^2$, on trouve que les seules fonctions \inserted{poly. de degré $\leq2$} quasi-invariantes sous $G$ sont \deleted{\unsure{les}} $a+bxy$ (qui sont \unsure{$=$|en|{}} invariantes)

Donc les seules côniques invariantes sous $G$ sont les coniques d'équation $a+bxy$ ; donc si, on veut les coniques circonscrites à $\Pi_\zeta$ il faut $a+b=0$ i.e. $b=-a$, on trouve la \deleted{\unsure{seule}} cônique \uline{invariante} circonscrite, d'équation $xy-1=0$ ....

% archive sheet number: 72
\scanpage{74}

En fait, le calcul en ``coordonnées épinglées universelles'' montre que c'est la seule cônique circonscrite (elle est donc automatiquement invariante ...)

\uline{Cas $n=3$} i.e. $\zeta=\zeta^{-2}$, $\zeta^{-1}=\zeta^{\unsure{2|1}}$, donc les fonctions pol. quasi-inv. \inserted{sous $G^+$ i.e. sous $u$} de degré $\leq2$ sont les

\[\left\{\begin{aligned}
a+bxy&&\text{val. propre }&1\\
ax+by^2&&\rule{1.6cm}{.35pt}&\zeta\\
a{\unsure{y|x}}^2+bx^2&&\rule{1.6cm}{.35pt}&\zeta^{-1}
\end{aligned}\right.\]

mais parmi ceux-ci, les qu. inv. sous $G$ sont encore $a+bxy$

\uline{Cas $n=4$} i.e. $\zeta^2=\zeta^{-2}$, les fonctions polyn. de degré $\leq2$ quasi-inv. sous $G^+$ \inserted{i.e. sous $u$} sont

\[\left\{\begin{aligned}
a+bxy&&\text{val. propre }&1\\
ax\quad(\text{resp }ay)&&\text{val. pr }&\zeta\quad(\zeta^{-1})\\
ax^2+by^2&&\text{val. propre }&\zeta^2=\zeta^{-2}
\end{aligned}\right.\]

Parmi celles-ci, les fonctions quasi-inv. sous $G$ tt entier sont les $a+bxy$, et \unsure{d'autre les|d'autre part les|distinctes} $ax^2+by^2$, \deleted{telles \unsure{telles} que}\inserted{ce qui \unsure{fait|font|forment} deux} familles à un paramètre de côniques inv. sous $G$. Dans chacune, il y a \deleted{\unsure{exactemt|exactement}} une cônique circonscrite à $S$, savoir l'habituelle $xy-1=0$, et de plus $x^2-y^2=0$

% archive sheet number: 73
